Video reasoning models increasingly act as general-purpose ``world models'' for downstream tasks such as robotics, simulation validation, and safety monitoring, yet existing video benchmarks often conflate temporal reasoning with static shortcuts, provide limited diagnostic structure, or emphasize generative plausibility rather than falsifiable physics checking.
We introduce \textbf{MIRAGE-V}, a taxonomy-driven benchmark for \emph{physics verification in videos} built from \emph{minimal-change paired clips} that differ by a single controlled intervention aligned to a specific failure category (e.g., temporal aliasing, occlusion-induced identity drift, latent physical property inference, and conservation-law violations).
Each example is annotated with a \emph{critical time window}, enabling sampling-robust evaluation and optional temporal localization.
We release an open-source generator for scalable pair creation, a unified evaluation harness supporting both open-weight and API-accessible models, and standardized analysis protocols with shortcut audits (text-only, single-frame, shuffled frames, and sparse-vs-dense sampling).
Across a \NumExamples-example evaluation spanning \NumModules\ taxonomy modules with \FramesPerVideo uniformly sampled frames per video, the best-performing model achieves \BestAcc\% accuracy but only \BestPairAcc\% pair accuracy, revealing that even frontier models frequently fail to discriminate minimally different clips and to reliably detect subtle physical inconsistencies.
These findings highlight a persistent gap between apparent video understanding and true physics checking, with important implications for deploying video reasoning systems in safety-critical and decision-making settings.
